\documentclass[10pt,a4paper]{article}
\usepackage{amsmath,amssymb,amsthm,graphicx}
\title{Exponential randomized-time lower bound for randomized Kt complexity and its gap version}
\author{Shuichi Hirahara}
\date{doi:10.4230/LIPIcs.CCC.2022.26; CC BY4.0}
\begin{document}
\maketitle
\noindent\textit{Editorial note.} Original author: Shuichi Hirahara. Official CCC2022 article26 (2022-07-11), DOI 10.4230/LIPIcs.CCC.2022.26. Creative Commons Attribution4.0.\par\medskip
\noindent\textit{Editorial note.} License: https://creativecommons.org/licenses/by/4.0/\par\medskip
\noindent\textit{Editorial note.} Faithful editable transcription of complete author Section5 Theorems5.3/5.8 and proofs, including exact prior reconstruction tool, author Corollary5.5/Lemma5.6 support and full page21 gap/i.o. continuation. Exact and gap variants count once. Original author mathematical expressions retained; no new proof or correction. Original CC BY4.0 license and complete bibliography retained.\par\medskip
\noindent\textit{Allowed standard prerequisites (editorial).}\begin{itemize}
\item Success-probability robustness of randomized Kt. Exact prior Lu-Oliveira Lemma5.2 retained.
\item Trevisan/Raz-Reingold-Vadhan reconstructive PRG construction. Author exact Lemma5.4 citing prior [32, Theorem4.7] and original Proof Sketch retained; this is a standard construction premise, not a replaced main proof.
\end{itemize}\medskip
\paragraph{3. Preliminaries: Notation.}
$[n]$ denotes $\{1,\ldots,n\}$. For a function $p:\mathbb N\to\mathbb N$ and $i\in\mathbb N$, the function $p^{(i)}:\mathbb N\to\mathbb N$ is recursively defined as follows: $p^{(0)}(n):=n$ and $p^{(i+1)}(n):=p^{(i)}(p(n))$ for every $n\in\mathbb N$.
\paragraph{Kolmogorov Complexity.}
For a string $x\in\{0,1\}^*$ and an index $i\in\mathbb N$, let $x_i$ denote the $i$-th bit of $x$ if $i\leq|x|$ and $\bot$ if $i>|x|$. We fix an efficient universal Turing machine $U$. Time-bounded Kolmogorov complexity is formally defined as follows.
\paragraph{Definition 3.1 (Time-bounded Kolmogorov complexity).}
For strings $x,y\in\{0,1\}^*$, a time bound $t\in\mathbb N\cup\{\infty\}$, and an oracle $A$, the $A$-oracle $t$-time-bounded Kolmogorov complexity of $x$ given $y$ is defined as
\[K^{t,A}(x\mid y):=\min\{|d|:U^{A,y,d}\text{ outputs }x_i\text{ on input }i\text{ in time }t\text{ for every }i\in[1,\ldots,|x|+1]\}.\]
Here, $U^{A,y,d}$ indicates that the universal Turing machine is given oracle access to $A$ and each bit of $y$ and $d\in\{0,1\}^*$. We omit the superscript $A$ if $A=\varnothing$, the superscript $t$ if $t=\infty$, and ``$\mid y$'' if $y$ is the empty string.
A simple counting argument implies the following basic fact of Kolmogorov-randomness.
\paragraph{Fact 3.2.}
For any integer $s\geq1$ and any string $y\in\{0,1\}^*$, the number of strings $x\in\{0,1\}^*$ such that $K(x\mid y)<s$ is less than $2^s$.
\begin{proof}
The number of programs of length less than $s$ is at most $\sum_{i=0}^{s-1}2^i<2^s$.
\end{proof}

\paragraph{5. A Lower Bound for Randomized Kt Complexity.}
Here, we implement Kolmogorov's insightful approach to the $\mathsf P$ versus $\mathsf{NP}$ problem for randomized $Kt$ complexity. We recall the formal definition of $rKt$, which was introduced by Oliveira [62].
\paragraph{Definition 5.1 ([62]).}
For every $x\in\{0,1\}^*$ and every $\lambda\in[0,1]$, the randomized time-bounded Kolmogorov complexity $rKt_\lambda(x)$ of $x$ is defined as
\[rKt_\lambda(x\mid y):=\min\left\{|d|+\lceil\log t\rceil:\Pr_{r\sim\{0,1\}^t}[U(d,y,r)\text{ outputs }x\text{ in time }t]\geq\lambda\right\}.\]
We omit the subscript ``$\lambda$'' if $\lambda=2/3$ and ``$\mid y$'' if $y$ is the empty string. The language $\mathsf{MrKtP}$ is defined as
\[\mathsf{MrKtP}:=\{(x,1^s):rKt(x)\leq s\}.\]
The definition of $rKt_\lambda$ is robust with respect to a choice of $\lambda$.
\paragraph{Lemma 5.2 (Lu and Oliveira [59]).}
For all $x\in\{0,1\}^*$ and $\lambda\in(0,1]$,
\[rKt(x)\leq rKt_\lambda(x)+O(\log(1/\lambda)).\]
In what follows, let $\epsilon>0$ be a sufficiently small positive constant. The $O$ notation below does not depend on $\epsilon$. We say that a pseudo-deterministic algorithm $M$ computes $y$ on input $x$ [27] if
\[\Pr_M[M(x)=y]\geq2/3,\]
where the probability is taken over an internal coin flip sequence of the randomized algorithm $M$. We use the following pseudorandom generator construction.
\paragraph{Lemma 5.4 (see the proof of [32, Theorem 4.7]).}
For all sufficiently large $n,m\in\mathbb N$ such that $m\leq2n$, there exists a triple $(G,A,R^{(-)})$ such that
\[\begin{aligned}
G&:\{0,1\}^n\times\{0,1\}^d\to\{0,1\}^m,\\
A&:\{0,1\}^n\times\{0,1\}^d\to\{0,1\}^m,\\
R^{(-)}&:\{0,1\}^m\times\{0,1\}^d\times\{0,1\}^r\to(\{0,1\}^n)^L,
\end{aligned}\]
and for every $x\in\{0,1\}^n$ and any $D:\{0,1\}^m\to\{0,1\}$ such that
\begin{multline*}
\left|\Pr_{\substack{z\sim\{0,1\}^d\\w'\sim\{0,1\}^{O(m)}}}[D(G(x;z);w')=1]\right.\\
\left.-\Pr_{\substack{w\sim\{0,1\}^m\\w'\sim\{0,1\}^{O(m)}}}[D(w;w')=1]\right|\geq\frac1m,
\end{multline*}
it holds that
\[\Pr_{\substack{w\sim\{0,1\}^r\\z\sim\{0,1\}^d}}[x\in R^D(A(x,z),z,w)]\geq\frac1{2m^2}.\]
Here, we have $d=O(\log^3n)$, $r=O(m)$, and $L=\operatorname{poly}(m)$. Moreover, $G$ and $A$ can be computed in time $\operatorname{poly}(n)$ and $R^D$ can be computed in time $\operatorname{poly}(n)$ with oracle access to $D$.
\begin{proof}[Proof Sketch]
The construction is given by the improvement of Trevisan's extractor [70] given by Raz et al. [64]. Specifically, we encode $x$ using a list-decodable error-correcting code and regard the encoding of $x$ as the truth table of a hard function $f$. Then, the pseudorandom generator construction $G$ is defined as the Nisan--Wigderson generator [61] instantiated with the function $f$ and the weak design of [64]. ([32, Theorem 4.7] is stated as a black-box hitting set generator construction; however, it is easy to observe that the construction actually provides a pseudorandom generator construction, which is a stronger object.)
\end{proof}

\paragraph{Theorem 5.3.}
There exists a constant $\epsilon>0$ such that $\mathsf{MrKtP}\notin\mathsf{BPTIME}(2^{\epsilon n})$.
\paragraph{Theorem 5.8 (Theorem 1.4, restated).}
$\mathsf{GapMrKtP}\notin\mathsf{i.o.BPTIME}(2^{\epsilon n})$ for some constant $\epsilon>0$.

\paragraph{Corollary 5.5.}
Under the same hypothesis of Lemma 5.4, it holds that
\[rKt^D(x)\leq m+O(\log^3 n).\]
\begin{proof}
By an averaging argument, there exists a string $z\in\{0,1\}^d$ such that
\[\Pr_w[x\in R^D(A(x,z),z,w)]\geq\frac1{2m^2}.\]
Let $\alpha:=A(x,z)\in\{0,1\}^m$. The string $x$ can be described by the following pseudo-deterministic algorithm $M$: $M$ takes the advice string $\alpha\in\{0,1\}^m$ and $z\in\{0,1\}^d$ as input, picks $w\sim\{0,1\}^r$ randomly, and outputs a random element of $R^D(\alpha,z,w)$. Since the list size is at most $L=\operatorname{poly}(m)$, the probability that the output of $M$ is equal to $x$ is at least $\lambda:=\frac1{2m^2}\cdot\frac1L$. Therefore, we obtain $rKt_\lambda(x)\leq m+d+O(\log n)$. The result follows from Lemma 5.2.
\end{proof}
\paragraph{Lemma 5.6.}
If $\mathsf{MrKtP}\in\mathsf{BPTIME}(2^{\epsilon n})$, then for all sufficiently large $n\in\mathbb N$ and any strings $x\in\{0,1\}^n$ and $y\in\{0,1\}^n$, it holds that
\[rKt(x\mid y)+rKt(y)\leq rKt(x,y)+O(\epsilon n).\]
\begin{proof}
We may assume without loss of generality that $rKt(y)\gg\log^3 n$. Let $m$ and $m'$ be parameters chosen later. Let $G:\{0,1\}^n\times\{0,1\}^d\to\{0,1\}^m$ and $G':\{0,1\}^n\times\{0,1\}^{d'}\to\{0,1\}^{m'}$ be the pseudorandom generator construction of Lemma 5.4 with parameters $m$ and $m'$, respectively.
Fix $z\in\{0,1\}^d$ and $z'\in\{0,1\}^{d'}$. Since the strings $(G(x;z),G'(y;z'))$ can be described by $z,z'$, and a program that computes $(x,y)$, we obtain
\begin{multline*}
rKt(G(x;z),G'(y;z'))\leq rKt(x,y)+|z|+|z'|+O(\log n)\\
=rKt(x,y)+O(\log^3 n).\tag{13}
\end{multline*}
Pick $w\sim\{0,1\}^m$ and $z'\sim\{0,1\}^{d'}$ randomly. We claim that $rKt(w,G'(y;z'))\geq|w|+|w'|-O(\log n)$ with high probability. Toward a contradiction, assume that
\[\Pr_{w,z'}[rKt(w,G'(y;z'))<|w|+|w'|-O(\log n)]\geq\frac12.\]
By Fact 3.2, we have
\[\Pr_{w,w'}[rKt(w,w')\geq|w|+|w'|-O(\log n)]\geq1-o(1).\]
We define an oracle $D$ so that $D(w';w):=1$ if and only if $rKt(w,w')<|w|+|w'|-O(\log n)$. Then, the two inequalities above show that $D$ distinguishes the output distribution of $G'(y;-)$ from the uniform distribution. By Corollary 5.5, we obtain
\[rKt^D(y)\leq m'+O(\log^3 n).\]
Using the assumption that $\mathsf{MrKtP}\in\mathsf{BPTIME}(2^{\epsilon n})$, the oracle $D$ can be computed by a randomized algorithm running in time $2^{O(\epsilon m')}$. It follows that
\[rKt(y)\leq m'+O(\epsilon m')=(1+O(\epsilon))\cdot m'.\]
We choose $m'$ so that this inequality does not hold; that is,
\[m':=(1-O(\epsilon))\cdot rKt(y).\]
Then, we obtain a contradiction and conclude that
\[\Pr_{w,z'}[rKt(w,G'(y;z'))\geq m+m'-O(\log n)]\geq\frac12.\tag{14}\]
Define $D_y$ to be an oracle such that $D_y(w;z'):=1$ if and only if $rKt(w,G'(y;z'))\leq rKt(x,y)+O(\log^3n)$. By Equations (14) and (13), $D_y$ distinguishes the output distribution of $G(x;-)$ from the uniform distribution if
\[m+m'-O(\log n)\geq rKt(x,y)+O(\log^3 n).\]
We define $m:=rKt(x,y)-m'+O(\log^3n)$ so that this inequality holds. By Corollary 5.5, we obtain
\[rKt^{D_y}(x)\leq m+O(\log^3n).\]
Since the oracle $D_y$ can be computed in time $2^{O(\epsilon m)}$ given $y$ as hard-wired input, we conclude that
\begin{align*}
rKt(x\mid y)&\leq m+O(\epsilon m).\\
&\leq rKt(x,y)-m'+O(\epsilon n).\\
&\leq rKt(x,y)-(1-O(\epsilon))\cdot rKt(y)+O(\epsilon n).\\
&\leq rKt(x,y)-rKt(y)+O(\epsilon n).
\end{align*}
\end{proof}
\begin{proof}[Proof of Theorem 5.3]
The outline of the proof is as follows: First, we find a string $y\in\{0,1\}^n$ such that $rKt(y)\geq n$ by an exhaustive search over all strings $y\in\{0,1\}^n$. Next, we find a string $x\in\{0,1\}^n$ such that $rKt(x,y)\geq2n-O(\epsilon n)$ by an exhaustive search over all strings $x\in\{0,1\}^n$. The existence of such a string $x$ is guaranteed by Lemma 5.6. Finally, we observe that the pair $(x,y)$ can be pseudo-deterministically computed in time $2^{n+O(\epsilon n)}$, which contradicts the lower bound on $rKt(x,y)$. Details follow.
Using the assumption that $\mathsf{MrKtP}\in\mathsf{BPTIME}(2^{\epsilon n})$, let $M$ be the randomized algorithm that computes $rKt(y)$ on input $y$ of length $n$ in time $2^{O(\epsilon n)}$ with probability at least $1-2^{-2n}$. Fix $n\in\mathbb N$. Let $y_n$ be the lexicographically first string $y\in\{0,1\}^n$ such that $rKt(y)\geq n$. Note that the existence of such a string $y_n$ is guaranteed by Fact 3.2.
We claim that there exists a pseudo-deterministic algorithm $M_1$ that computes $y_n$ on input $n\in\mathbb N$ in time $2^{n+O(\epsilon n)}$. For all strings $y\in\{0,1\}^n$ in lexicographical order, the algorithm $M_1$ tests whether $rKt(y)\geq n$ using $M$, and outputs the first string $y$ that passes this test. Since the error probability of $M$ is at most $2^{-2n}$, by a union bound, the algorithm $M_1$ computes $y_n$ pseudo-deterministically with probability $1-2^{-n}$. The running time of $M_1$ is at most $2^{n+O(\epsilon n)}$.
Next, let $x_n$ be the lexicographically first string $x\in\{0,1\}^n$ such that
\[rKt(x,y_n)\geq2n-c\cdot\epsilon n,\tag{15}\]
where $c$ is a large constant to be chosen independently of $\epsilon$. We prove that $x_n$ is well defined, i.e., there exists a string $x\in\{0,1\}^n$ such that $rKt(x,y_n)\geq2n-c\cdot\epsilon n$. By Fact 3.2, there exists a string $x\in\{0,1\}^n$ such that $rKt(x\mid y)\geq n$. By Lemma 5.6,
\[rKt(x,y)+O(\epsilon n)\geq rKt(x\mid y)+rKt(y)\geq2n.\]
Choosing a large constant $c$, we obtain $rKt(x,y)\geq2n-c\cdot\epsilon n$, as desired.
We claim that there exists a pseudo-deterministic algorithm $M_2$ that computes $x_n$ given $y_n$ as input in time $2^{n+O(\epsilon n)}$. For all strings $x\in\{0,1\}^n$ in lexicographical order, the algorithm $M_2$ tests whether $rKt(x,y_n)\geq2n-c\cdot\epsilon n$ using $M$, and outputs the first string $x$ that passes this test. Using a union bound, the error probability of $M_2$ is bounded by $2^{-n}$. The running time of $M_2$ is at most $2^{n+O(\epsilon n)}$.
Finally, by combining the algorithms $M_1$ and $M_2$, we obtain an algorithm that computes $(x_n,y_n)$ pseudo-deterministically in time $2^{n+O(\epsilon n)}$. By the definition of $rKt$, this implies that
\[rKt(x_n,y_n)\leq n+O(\epsilon n).\]
We also have $rKt(x_n,y_n)\geq2n-O(\epsilon n)$ by Equation (15). Therefore, we obtain $2n-O(\epsilon n)\leq n+O(\epsilon n)$, which is a contradiction for a sufficiently small constant $\epsilon>0$.
\end{proof}
We observe that $\mathsf{MrKtP}$ is in the exponential-time variant of $\mathsf{PP}$.
\paragraph{Proposition 5.7.}
$\mathsf{MrKtP}\in\mathsf{PEXP}$.
\begin{proof}
Since $\mathsf{PP}$ is closed under truth-table reductions [26], it suffices to show that $\mathsf{MrKtP}$ is reducible to $\mathsf{PP}$ via an exponential-time truth-table reduction. Let $(x,1^s)$ be an input such that $x\in\{0,1\}^n$ and $s\in\mathbb N$. For each string $d\in\{0,1\}^*$ of length at most $s$, we ask the $\mathsf{PP}$ oracle whether $U(d,r)$ outputs $x$ in time $t$ for at least a $2/3$-fraction of $r\in\{0,1\}^t$, where $t:=2^{s-|d|}$. We accept the input $(x,1^s)$ if and only if a positive answer is returned from the oracle.
\end{proof}
We expect that Theorem 5.3 can be extended to a lower bound against $\mathsf{PP}$-type algorithms. However, the original motivation of Oliveira [62] is to study a gap version of $\mathsf{MrKtP}$, for which the upper bound of $\mathsf{pr\text{-}BPE}$ can be proved. Specifically, let $\mathsf{GapMrKtP}$ denote the promise problem $(\Pi_{\mathrm{Yes}},\Pi_{\mathrm{No}})$ such that $\Pi_{\mathrm{Yes}}=\{(x,s):rKt(x)\leq s\}$ and $\Pi_{\mathrm{No}}=\{(x,s):rKt(x)>s+c\log|x|\}$, where $c$ is a sufficiently large constant. One can observe $\mathsf{GapMrKtP}\in\mathsf{pr\text{-}BPE}$ [62]. Theorem 5.3 can be extended to this promise problem.
Here, abusing a notation, for a class $C$ of languages, we say that a promise problem $\Pi\in C$ if there exists a language $L\in C$ such that $\Pi_{\mathrm{Yes}}\subseteq L\subseteq\{0,1\}^*\setminus\Pi_{\mathrm{No}}$ for $(\Pi_{\mathrm{Yes}},\Pi_{\mathrm{No}})=\Pi$.
The lower bound of Theorem 5.8 does work against a bounded probabilistic algorithm that satisfies the promise of $\mathsf{BPP}$-type algorithms over all the inputs. In contrast, the upper bound $\mathsf{pr\text{-}BPE}=\mathsf{pr\text{-}BPTIME}(2^{O(n)})$ holds only for an algorithm that may not satisfy the promise of $\mathsf{BPP}$-type algorithms for some inputs. This leaves a qualitative gap between the lower bound and the upper bound on $\mathsf{GapMrKtP}$; closing this gap is an interesting open question. We note, however, that there is no gap under the plausible assumption that $\mathsf E\not\subseteq\mathsf{i.o.SIZE}(2^{\epsilon n})$ for some constant $\epsilon>0$, in which case $\mathsf{pr\text{-}BPP}=\mathsf P=\mathsf{BPP}$ [46].
\begin{proof}[Proof of Theorem 5.8]
The proof is essentially the same with Theorem 5.3. We only explain how to modify the proof of Theorem 5.3.
Toward a contradiction, we assume the existence of a randomized algorithm $M$ that computes some language $L$ infinitely often, where $L$ is a language consistent with $\mathsf{GapMrKtP}$. For all large $n\in\mathbb N$ and every $x\in\{0,1\}^n$ and $s\in[n]$, we assume that the binary encoding of $(x,s)$ is exactly equal to $m(n)$, where $m(n)=\Theta(n)$ is some function. This ensures that for infinitely many $n\in\mathbb N$, the algorithm $M$ computes $L$ correctly on input $(x,s)$ for every $x\in\{0,1\}^n$ and every $s\in[n]$.
First, we explain how to extend the lower bound of Theorem 5.3 to the infinitely-often version. Let $n$ be the length of $x_n\in\{0,1\}^n$ and $y_n\in\{0,1\}^n$ constructed in the proof of Theorem 5.3. By inspection, in the proof of Theorem 5.3, we use the algorithm $M$ on inputs of length at most $n'$, where $n'=O(n)$. Any input $x$ of length less than $n'$ can be padded to an input of length $n'$ by mapping $x$ to $x':=x1^{n'-|x|}$. This ensures that the length of any input to $M$ that we use in the proof is exactly equal to $n'$. Moreover, the correctness of the proof remains unchanged because $rKt(x)=rKt(x')\pm O(\log n)$.
Second, we explain how to deal with the promise problem. In the proof of Theorem 5.3, we exhaustively search the lexicographically first string $y_n\in\{0,1\}^n$ such that $rKt(y_n)\geq n$. Instead, we search the lexicographically first string $y_n\in\{0,1\}^n$ such that $(y_n,n-2c\cdot\log n)\notin L$. Since $L\subseteq\{0,1\}^*\setminus\Pi_{\mathrm{No}}$, this ensures that $rKt(y_n)\geq n-c\log n$. Similarly, we search the lexicographically first string $x_n\in\{0,1\}^n$ such that $((x_n,y_n),2n-O(\epsilon n))\notin L$. This ensures that $rKt(x_n,y_n)\geq2n-O(\epsilon n)$. Since the language $L$ can be decided by the randomized algorithm $M$, we can compute the pair $(x_n,y_n)$ pseudo-deterministically in time $2^{n+O(\epsilon n)}$. This is a contradiction.
\end{proof}

\begin{thebibliography}{99}
\bibitem[1]{published-1} Dorit Aharonov and Oded Regev. Lattice problems in NP $\cap$ coNP. J. ACM, 52(5):749–765, 2005. doi:10.1145/1089023.1089025.
\bibitem[2]{published-2} Miklós Ajtai. Generating Hard Instances of Lattice Problems (Extended Abstract). In Proceedings of the Symposium on the Theory of Computing (STOC), pages 99–108, 1996. doi:10.1145/237814.237838.
\bibitem[3]{published-3} Eric Allender, Harry Buhrman, Michal Koucký, Dieter van Melkebeek, and Detlef Ronneburger. Power from Random Strings. SIAM J. Comput., 35(6):1467–1493, 2006. doi:10.1137/ 050628994.
\bibitem[4]{published-4} Eric Allender, Mahdi Cheraghchi, Dimitrios Myrisiotis, Harsha Tirumala, and Ilya Volkovich. One-Way Functions and a Conditional Variant of MKTP. In Proceedings of the Foundations of Software Technology and Theoretical Computer Science (FSTTCS), pages 7:1–7:19, 2021. doi:10.4230/LIPIcs.FSTTCS.2021.7.
\bibitem[5]{published-5} Eric Allender, John Gouwar, Shuichi Hirahara, and Caleb Robelle. Cryptographic Hard- ness Under Projections for Time-Bounded Kolmogorov Complexity. In Proceedings of the International Symposium on Algorithms and Computation (ISAAC), pages 54:1–54:17, 2021. doi:10.4230/LIPIcs.ISAAC.2021.54.
\bibitem[6]{published-6} Eric Allender, Joshua A. Grochow, Dieter van Melkebeek, Cristopher Moore, and Andrew Morgan. Minimum Circuit Size, Graph Isomorphism, and Related Problems. SIAM J. Comput., 47(4):1339–1372, 2018. doi:10.1137/17M1157970.
\bibitem[7]{published-7} Eric Allender, Lisa Hellerstein, Paul McCabe, Toniann Pitassi, and Michael E. Saks. Minimizing Disjunctive Normal Form Formulas and AC0 Circuits Given a Truth Table. SIAM J. Comput., 38(1):63–84, 2008. doi:10.1137/060664537.
\bibitem[8]{published-8} Eric Allender, Michal Koucký, Detlef Ronneburger, and Sambuddha Roy. The pervasive reach of resource-bounded Kolmogorov complexity in computational complexity theory. J. Comput. Syst. Sci., 77(1):14–40, 2011. doi:10.1016/j.jcss.2010.06.004.
\bibitem[9]{published-9} Luis Antunes, Lance Fortnow, Dieter van Melkebeek, and N. V. Vinodchandran. Computational depth: Concept and applications. Theor. Comput. Sci., 354(3):391–404, 2006. doi:10.1016/ j.tcs.2005.11.033.
\bibitem[10]{published-10} Luis Antunes, Sophie Laplante, Alexandre Pinto, and Liliana C. M. Salvador. Cryptographic Security of Individual Instances. In Proceedings of Information Theoretic Security (ICITS), pages 195–210, 2007. doi:10.1007/978-3-642-10230-1\_17.
\bibitem[11]{published-11} Luis Filipe Coelho Antunes and Lance Fortnow. Worst-Case Running Times for Average-Case Algorithms. In Proceedings of the Conference on Computational Complexity (CCC), pages 298–303, 2009. doi:10.1109/CCC.2009.12.
\bibitem[12]{published-12} Amos Beimel. Secret-Sharing Schemes: A Survey. In The Third International Workshop on Coding and Cryptology (IWCC), pages 11–46, 2011. doi:10.1007/978-3-642-20901-7\_2.
\bibitem[13]{published-13} Josh Cohen Benaloh and Jerry Leichter. Generalized Secret Sharing and Monotone Functions. In Proceedings of the International Cryptology Conference (CRYPTO), pages 27–35, 1988. doi:10.1007/0-387-34799-2\_3.
\bibitem[14]{published-14} George Robert Blakley. Safeguarding cryptographic keys. In International Workshop on Managing Requirements Knowledge (MARK), pages 313–318, 1979. doi:10.1109/MARK.1979. 8817296.
\bibitem[15]{published-15} Manuel Blum and Sampath Kannan. Designing Programs that Check Their Work. J. ACM, 42(1):269–291, 1995. doi:10.1145/200836.200880.
\bibitem[16]{published-16} Andrej Bogdanov and Luca Trevisan. Average-Case Complexity. Foundations and Trends in Theoretical Computer Science, 2(1), 2006. doi:10.1561/0400000004.
\bibitem[17]{published-17} Andrej Bogdanov and Luca Trevisan. On Worst-Case to Average-Case Reductions for NP Problems. SIAM J. Comput., 36(4):1119–1159, 2006. doi:10.1137/S0097539705446974.
\bibitem[18]{published-18} Harry Buhrman, Lance Fortnow, and Aduri Pavan. Some Results on Derandomization. Theory Comput. Syst., 38(2):211–227, 2005. doi:10.1007/s00224-004-1194-y.
\bibitem[19]{published-19} Harry Buhrman and Elvira Mayordomo. An Excursion to the Kolmogorov Random Strings. J. Comput. Syst. Sci., 54(3):393–399, 1997. doi:10.1006/jcss.1997.1484.
\bibitem[20]{published-20} Lijie Chen, Shuichi Hirahara, and Neekon Vafa. Average-case Hardness of NP and PH from Worst-case Fine-grained Assumptions. In Proceedings of the Innovations in Theoretical Computer Science Conference (ITCS), pages 67:1–67:17, 2022.
\bibitem[21]{published-21} Irit Dinur, Prahladh Harsha, and Guy Kindler. Polynomially Low Error PCPs with polyloglog n Queries via Modular Composition. In Proceedings of the Symposium on Theory of Computing (STOC), pages 267–276, 2015. doi:10.1145/2746539.2746630.
\bibitem[22]{published-22} Irit Dinur and Shmuel Safra. On the hardness of approximating label-cover. Inf. Process. Lett., 89(5):247–254, 2004. doi:10.1016/j.ipl.2003.11.007.
\bibitem[23]{published-23} Joan Feigenbaum and Lance Fortnow. Random-Self-Reducibility of Complete Sets. SIAM J. Comput., 22(5):994–1005, 1993. doi:10.1137/0222061.
\bibitem[24]{published-24} Lance Fortnow and Martin Kummer. On Resource-Bounded Instance Complexity. Theor. Comput. Sci., 161(1\&2):123–140, 1996. doi:10.1016/0304-3975(95)00097-6.
\bibitem[25]{published-25} Lance Fortnow, Troy Lee, and Nikolai K. Vereshchagin. Kolmogorov Complexity with Error. In STACS 2006, 23rd Annual Symposium on Theoretical Aspects of Computer Science, Marseille, France, February 23-25, 2006, Proceedings, pages 137–148, 2006. doi:10.1007/11672142\_10.
\bibitem[26]{published-26} Lance Fortnow and Nick Reingold. PP is Closed Under Truth-Table Reductions. Inf. Comput., 124(1):1–6, 1996. doi:10.1006/inco.1996.0001.
\bibitem[27]{published-27} Eran Gat and Shafi Goldwasser. Probabilistic Search Algorithms with Unique Answers and Their Cryptographic Applications. Electron. Colloquium Comput. Complex., page 136, 2011. URL: https://eccc.weizmann.ac.il/report/2011/136.
\bibitem[28]{published-28} Halley Goldberg and Valentine Kabanets. A Simpler Proof of the Worst-Case to Average-Case Reduction for Polynomial Hierarchy via Symmetry of Information. Electronic Colloquium on Computational Complexity (ECCC), 007, 2022.
\bibitem[29]{published-29} Oded Goldreich and Shafi Goldwasser. On the Limits of Nonapproximability of Lattice Problems. J. Comput. Syst. Sci., 60(3):540–563, 2000. doi:10.1006/jcss.1999.1686.
\bibitem[30]{published-30} Ishay Haviv and Oded Regev. Tensor-based Hardness of the Shortest Vector Problem to within Almost Polynomial Factors. Theory of Computing, 8(1):513–531, 2012. doi:10.4086/ toc.2012.v008a023.
\bibitem[31]{published-31} Shuichi Hirahara. Non-black-box Worst-case to Average-case Reductions within NP. In Proceedings of the Symposium on Foundations of Computer Science (FOCS), pages 247–258, 2018.
\bibitem[32]{published-32} Shuichi Hirahara. Characterizing Average-Case Complexity of PH by Worst-Case Meta- Complexity. In Proceedings of the Symposium on Foundations of Computer Science (FOCS), pages 50–60, 2020.
\bibitem[33]{published-33} Shuichi Hirahara. Non-Disjoint Promise Problems from Meta-Computational View of Pseudor- andom Generator Constructions. In Proceedings of the Computational Complexity Conference (CCC), pages 20:1–20:47, 2020. doi:10.4230/LIPIcs.CCC.2020.20.
\bibitem[34]{published-34} Shuichi Hirahara. Unexpected hardness results for Kolmogorov complexity under uniform reductions. In Proceedings of the Symposium on Theory of Computing (STOC), pages 1038– 1051, 2020. doi:10.1145/3357713.3384251.
\bibitem[35]{published-35} Shuichi Hirahara. Average-case hardness of NP from exponential worst-case hardness assump- tions. In Proceedings of the Symposium on Theory of Computing (STOC), pages 292–302, 2021. doi:10.1145/3406325.3451065.
\bibitem[36]{published-36} Shuichi Hirahara and Mikito Nanashima. On Worst-Case Learning in Relativized Heuristica. In Proceedings of the Symposium on Foundations of Computer Science (FOCS), 2021.
\bibitem[37]{published-37} Shuichi Hirahara, Igor Carboni Oliveira, and Rahul Santhanam. NP-hardness of Minimum Circuit Size Problem for OR-AND-MOD Circuits. In Proceedings of the Computational Complexity Conference (CCC), pages 5:1–5:31, 2018. doi:10.4230/LIPIcs.CCC.2018.5.
\bibitem[38]{published-38} Shuichi Hirahara and Rahul Santhanam. Errorless versus Error-prone Average-case Complexity. In Proceedings of the Innovations in Theoretical Computer Science Conference (ITCS), pages 38:1–38:23, 2022.
\bibitem[39]{published-39} Shuichi Hirahara and Osamu Watanabe. Limits of Minimum Circuit Size Problem as Oracle. In Proceedings of the Conference on Computational Complexity (CCC), pages 18:1–18:20, 2016. doi:10.4230/LIPIcs.CCC.2016.18.
\bibitem[40]{published-40} Rahul Ilango. Approaching MCSP from Above and Below: Hardness for a Conditional Variant and AC0 [p]. In Proceedings of the Innovations in Theoretical Computer Science Conference (ITCS), pages 34:1–34:26, 2020. doi:10.4230/LIPIcs.ITCS.2020.34.
\bibitem[41]{published-41} Rahul Ilango. Constant Depth Formula and Partial Function Versions of MCSP are Hard. In Proceedings of the Symposium on Foundations of Computer Science (FOCS), pages 424–433, 2020.
\bibitem[42]{published-42} Rahul Ilango, Bruno Loff, and Igor Carboni Oliveira. NP-Hardness of Circuit Minimization for Multi-Output Functions. In Proceedings of the Computational Complexity Conference (CCC), pages 22:1–22:36, 2020. doi:10.4230/LIPIcs.CCC.2020.22.
\bibitem[43]{published-43} Russell Impagliazzo. A Personal View of Average-Case Complexity. In Proceedings of the Structure in Complexity Theory Conference, pages 134–147, 1995. doi:10.1109/SCT.1995. 514853.
\bibitem[44]{published-44} Russell Impagliazzo. Relativized Separations of Worst-Case and Average-Case Complexities for NP. In Proceedings of the Conference on Computational Complexity (CCC), pages 104–114, 2011. doi:10.1109/CCC.2011.34.
\bibitem[45]{published-45} Russell Impagliazzo and Leonid A. Levin. No Better Ways to Generate Hard NP Instances than Picking Uniformly at Random. In Proceedings of the Symposium on Foundations of Computer Science (FOCS), pages 812–821, 1990. doi:10.1109/FSCS.1990.89604.
\bibitem[46]{published-46} Russell Impagliazzo and Avi Wigderson. P = BPP if E Requires Exponential Circuits: Derandomizing the XOR Lemma. In Proceedings of the Symposium on the Theory of Computing (STOC), pages 220–229, 1997. doi:10.1145/258533.258590.
\bibitem[47]{published-47} Valentine Kabanets and Jin-yi Cai. Circuit minimization problem. In Proceedings of the Symposium on Theory of Computing (STOC), pages 73–79, 2000. doi:10.1145/335305. 335314.
\bibitem[48]{published-48} Tarik Kaced. Almost-perfect secret sharing. In Proceedings of the International Symposium on Information Theory (ISIT), pages 1603–1607, 2011. doi:10.1109/ISIT.2011.6033816.
\bibitem[49]{published-49} Mauricio Karchmer and Avi Wigderson. On Span Programs. In Proceedings of the Structure in Complexity Theory Conference, pages 102–111, 1993. doi:10.1109/SCT.1993.336536.
\bibitem[50]{published-50} Ker-I Ko. On the Complexity of Learning Minimum Time-Bounded Turing Machines. SIAM J. Comput., 20(5):962–986, 1991. doi:10.1137/0220059.
\bibitem[51]{published-51} Andrei N Kolmogorov. Three approaches to the quantitative definition of information. Problems of information transmission, 1(1):1–7, 1965.
\bibitem[52]{published-52} Troy Lee and Andrei E. Romashchenko. Resource bounded symmetry of information revisited. Theor. Comput. Sci., 345(2-3):386–405, 2005. doi:10.1016/j.tcs.2005.07.017.
\bibitem[53]{published-53} Leonid A. Levin. The Tale of One-Way Functions. Probl. Inf. Transm., 39(1):92–103, 2003. doi:10.1023/A:1023634616182.
\bibitem[54]{published-54} Leonid A. Levin. How do we succeed in tasks like proving Fermat’s Theorem or predicting the Higgs boson?, 2021. An invited talk given at STOC 2021. The transcript is available at https://www.cs.bu.edu/fac/lnd/expo/stoc21/txt.pdf.
\bibitem[55]{published-55} Ming Li and Paul M. B. Vitányi. An Introduction to Kolmogorov Complexity and Its Applications, Third Edition. Texts in Computer Science. Springer, 2008. doi:10.1007/ 978-0-387-49820-1.
\bibitem[56]{published-56} Yanyi Liu and Rafael Pass. On One-way Functions from NP-Complete Problems. Electron. Colloquium Comput. Complex., 28:59, 2021. URL: https://eccc.weizmann.ac.il/report/ 2021/059.
\bibitem[57]{published-57} Luc Longpré and Sarah Mocas. Symmetry of Information and One-Way Functions. Inf. Process. Lett., 46(2):95–100, 1993. doi:10.1016/0020-0190(93)90204-M.
\bibitem[58]{published-58} Luc Longpré and Osamu Watanabe. On Symmetry of Information and Polynomial Time Invertibility. Inf. Comput., 121(1):14–22, 1995. doi:10.1006/inco.1995.1120.
\bibitem[59]{published-59} Zhenjian Lu and Igor Carboni Oliveira. An Efficient Coding Theorem via Probabilistic Representations and Its Applications. In Proceedings of the International Colloquium on Automata, Languages, and Programming (ICALP), pages 94:1–94:20, 2021. doi:10.4230/ LIPIcs.ICALP.2021.94.
\bibitem[60]{published-60} Daniele Micciancio and Oded Regev. Worst-Case to Average-Case Reductions Based on Gaus- sian Measures. SIAM J. Comput., 37(1):267–302, 2007. doi:10.1137/S0097539705447360.
\bibitem[61]{published-61} Noam Nisan and Avi Wigderson. Hardness vs Randomness. J. Comput. Syst. Sci., 49(2):149– 167, 1994. doi:10.1016/S0022-0000(05)80043-1.
\bibitem[62]{published-62} Igor Carboni Oliveira. Randomness and Intractability in Kolmogorov Complexity. In Proceed- ings of the International Colloquium on Automata, Languages, and Programming (ICALP), pages 32:1–32:14, 2019. doi:10.4230/LIPIcs.ICALP.2019.32.
\bibitem[63]{published-63} Igor Carboni Oliveira and Rahul Santhanam. Pseudodeterministic constructions in subex- ponential time. In Proceedings of the Symposium on Theory of Computing (STOC), pages 665–677, 2017. doi:10.1145/3055399.3055500.
\bibitem[64]{published-64} Ran Raz, Omer Reingold, and Salil P. Vadhan. Extracting all the Randomness and Reducing the Error in Trevisan’s Extractors. J. Comput. Syst. Sci., 65(1):97–128, 2002. doi:10.1006/ jcss.2002.1824.
\bibitem[65]{published-65} Hanlin Ren and Rahul Santhanam. Hardness of KT Characterizes Parallel Cryptography. In Proceedings of the Computational Complexity Conference (CCC), pages 35:1–35:58, 2021. doi:10.4230/LIPIcs.CCC.2021.35.
\bibitem[66]{published-66} Detlef Ronneburger. Kolmogorov Complexity and Derandomization. PhD thesis, Rutgers UniversityDept. of Computer Science Laboratory for Computer Sci. Research Hill Center, NJ, USA, 2004.
\bibitem[67]{published-67} Adi Shamir. How to Share a Secret. Commun. ACM, 22(11):612–613, 1979. doi:10.1145/ 359168.359176.
\bibitem[68]{published-68} Amnon Ta-Shma, Christopher Umans, and David Zuckerman. Lossless Condensers, Un- balanced Expanders, And Extractors. Combinatorica, 27(2):213–240, 2007. doi:10.1007/ s00493-007-0053-2.
\bibitem[69]{published-69} Boris A. Trakhtenbrot. A Survey of Russian Approaches to Perebor (Brute-Force Searches) Algorithms. IEEE Annals of the History of Computing, 6(4):384–400, 1984. doi:10.1109/ MAHC.1984.10036.
\bibitem[70]{published-70} Luca Trevisan. Extractors and pseudorandom generators. J. ACM, 48(4):860–879, 2001. doi:10.1145/502090.502099.
\bibitem[71]{published-71} Luca Trevisan and Salil P. Vadhan. Pseudorandomness and Average-Case Complexity Via Uniform Reductions. Computational Complexity, 16(4):331–364, 2007. doi:10.1007/ s00037-007-0233-x.
\bibitem[72]{published-72} Emanuele Viola. On Constructing Parallel Pseudorandom Generators from One-Way Functions. In Proceedings of the Conference on Computational Complexity (CCC), pages 183–197, 2005. doi:10.1109/CCC.2005.16.
\bibitem[73]{published-73} Emanuele Viola. The complexity of constructing pseudorandom generators from hard functions. Computational Complexity, 13(3-4):147–188, 2005. doi:10.1007/s00037-004-0187-1.
\bibitem[74]{published-74} AK Zvonkin and LA Levin. The complexity of finite objects and the algorithmic concepts of randomness and information. UMN (Russian Math. Surveys), 25(6):83–124, 1970.
\end{thebibliography}
\end{document}
